\documentclass[11pt]{article}
\usepackage[margin=1in]{geometry}
\usepackage{amsmath, amssymb, amsthm}
\usepackage{algorithm}
\usepackage{algpseudocode}
\usepackage{booktabs}
\usepackage{hyperref}

\newtheorem{theorem}{Theorem}
\newtheorem{lemma}[theorem]{Lemma}
\newtheorem{definition}{Definition}
\newtheorem{property}{Property}

\title{Mathematical Specification and Variance Scaling Theory of Linear Affine Layers (ALGO-NN-03)}
\author{Algorithms Engineering Group}
\date{\today}

\begin{document}

\maketitle

\begin{abstract}
This document formulates the mathematical properties, algebraic representations, and variance preservation theory of linear (affine) transformations in deep neural networks. We derive the exact fan-in and fan-out variance scaling conditions under Glorot (Xavier), He (Kaiming), and LeCun initializations, prove variance conservation across forward and backward propagation steps, establish computational complexity bounds, and provide a fully computed numeric trace matching the production implementation contract.
\end{abstract}

\section{Notation}
\label{sec:notation}

Let $\mathbb{R}$ denote the set of real numbers and $\mathbb{N} = \{1, 2, \dots\}$. Matrices are denoted by uppercase bold letters ($\mathbf{X}, \mathbf{W}, \mathbf{Y}$), vectors by lowercase bold letters ($\mathbf{x}, \mathbf{b}, \mathbf{y}$), and scalars by standard italics.

\begin{itemize}
    \item $B \in \mathbb{N}$: Batch size (number of independent examples).
    \item $d_{\text{in}} \in \mathbb{N}$: Fan-in dimension (input feature dimensionality).
    \item $d_{\text{out}} \in \mathbb{N}$: Fan-out dimension (output feature dimensionality).
    \item $\mathbf{X} \in \mathbb{R}^{B \times d_{\text{in}}}$: Input batch matrix, where row $i$ represents sample $\mathbf{x}_i \in \mathbb{R}^{d_{\text{in}}}$.
    \item $\mathbf{W} \in \mathbb{R}^{d_{\text{out}} \times d_{\text{in}}}$: Parameterized weight matrix.
    \item $\mathbf{b} \in \mathbb{R}^{d_{\text{out}}}$: Parameterized bias vector.
    \item $\mathbf{Y} \in \mathbb{R}^{B \times d_{\text{out}}}$: Output batch matrix.
    \item $\mathbb{E}[\cdot]$: Expectation operator.
    \item $\text{Var}(\cdot)$: Variance operator.
\end{itemize}

\section{Problem Definition}
\label{sec:problem_definition}

\begin{definition}[Affine Transformation]
Given input matrix $\mathbf{X} \in \mathbb{R}^{B \times d_{\text{in}}}$, weight matrix $\mathbf{W} \in \mathbb{R}^{d_{\text{out}} \times d_{\text{in}}}$, and bias vector $\mathbf{b} \in \mathbb{R}^{d_{\text{out}}}$, the affine layer computes the linear operator:
\begin{equation}
\label{eq:affine_layer}
\mathbf{Y} = \mathbf{X} \mathbf{W}^\top + \mathbf{1}_B \mathbf{b}^\top,
\end{equation}
where $\mathbf{1}_B = [1, 1, \dots, 1]^\top \in \mathbb{R}^B$. For each individual element $(i, j)$:
\begin{equation}
Y_{ij} = \sum_{k=1}^{d_{\text{in}}} X_{ik} W_{jk} + b_j, \quad \forall i \in \{1, \dots, B\}, \; j \in \{1, \dots, d_{\text{out}}\}.
\end{equation}
\end{definition}

\begin{definition}[Fan-in and Fan-out]
For a linear transformation with weight matrix $\mathbf{W} \in \mathbb{R}^{d_{\text{out}} \times d_{\text{in}}}$:
\begin{align}
\text{fan\_in} &= d_{\text{in}} = \text{number of columns in } \mathbf{W}, \\
\text{fan\_out} &= d_{\text{out}} = \text{number of rows in } \mathbf{W}.
\end{align}
\end{definition}

\section{Algorithm}
\label{sec:algorithm}

Algorithm~\ref{alg:linear_layer} specifies the deterministic forward execution and variance metric calculation.

\begin{algorithm}[H]
\caption{Affine Matrix Multiplication and Variance Scaling Calculation}
\label{alg:linear_layer}
\begin{algorithmic}[1]
\Require Input batch $\mathbf{X} \in \mathbb{R}^{B \times d_{\text{in}}}$, weights $\mathbf{W} \in \mathbb{R}^{d_{\text{out}} \times d_{\text{in}}}$, bias $\mathbf{b} \in \mathbb{R}^{d_{\text{out}}}$, mode $\in \{\text{kaiming}, \text{glorot}, \text{lecun}\}$.
\Ensure Transformed matrix $\mathbf{Y} \in \mathbb{R}^{B \times d_{\text{out}}}$, fan metrics, theoretical variance $\sigma^2$, bound $r$.
\State $B \gets \text{rows}(\mathbf{X})$, $d_{\text{in}} \gets \text{cols}(\mathbf{X})$, $d_{\text{out}} \gets \text{rows}(\mathbf{W})$
\State Initialize $\mathbf{Y} \in \mathbb{R}^{B \times d_{\text{out}}}$ to zeros
\For{$i = 1 \textbf{ to } B$}
    \For{$j = 1 \textbf{ to } d_{\text{out}}$}
        \State $\text{sum} \gets 0.0$
        \For{$k = 1 \textbf{ to } d_{\text{in}}$}
            \State $\text{sum} \gets \text{sum} + X_{ik} \cdot W_{jk}$
        \EndFor
        \State $Y_{ij} \gets \text{sum} + b_j$
    \EndFor
\EndFor
\If{$\text{mode} = \text{glorot\_uniform}$}
    \State $\sigma^2 \gets \frac{2}{d_{\text{in}} + d_{\text{out}}}$, $r \gets \sqrt{3 \sigma^2} = \sqrt{\frac{6}{d_{\text{in}} + d_{\text{out}}}}$
\ElsIf{$\text{mode} = \text{kaiming\_uniform}$}
    \State $\sigma^2 \gets \frac{2}{d_{\text{in}}}$, $r \gets \sqrt{3 \sigma^2} = \sqrt{\frac{6}{d_{\text{in}}}}$
\ElsIf{$\text{mode} = \text{lecun\_uniform}$}
    \State $\sigma^2 \gets \frac{1}{d_{\text{in}}}$, $r \gets \sqrt{3 \sigma^2} = \sqrt{\frac{3}{d_{\text{in}}}}$
\EndIf
\State \Return $(\mathbf{Y}, d_{\text{in}}, d_{\text{out}}, B, \sigma^2, r)$
\end{algorithmic}
\end{algorithm}

\section{Variance Scaling and Initialization Theory}
\label{sec:correctness}

To prevent activations and backpropagated gradients from exponentially vanishing or exploding across deep networks, the variance of the signal must be preserved layer-to-layer.

\begin{theorem}[Variance Preservation in Linear Maps]
\label{thm:variance_preservation}
Let $y_j = \sum_{k=1}^{d_{\text{in}}} W_{jk} x_k$ with bias $b_j = 0$. Assume that $x_k$ and $W_{jk}$ are independent, identically distributed random variables with zero mean ($\mathbb{E}[x_k] = 0, \mathbb{E}[W_{jk}] = 0$) and finite variances $\text{Var}(x)$ and $\text{Var}(W)$. Then:
\begin{equation}
\label{eq:var_forward}
\text{Var}(y_j) = d_{\text{in}} \cdot \text{Var}(W) \cdot \text{Var}(x).
\end{equation}
\end{theorem}

\begin{proof}
Because $W_{jk}$ and $x_k$ are independent with zero mean:
\begin{align}
\mathbb{E}[y_j] &= \sum_{k=1}^{d_{\text{in}}} \mathbb{E}[W_{jk}] \mathbb{E}[x_k] = 0, \\
\text{Var}(y_j) &= \mathbb{E}[y_j^2] = \mathbb{E}\left[ \left(\sum_{k=1}^{d_{\text{in}}} W_{jk} x_k\right)^2 \right] = \sum_{k=1}^{d_{\text{in}}} \sum_{m=1}^{d_{\text{in}}} \mathbb{E}[W_{jk} W_{jm} x_k x_m].
\end{align}
Since cross terms ($k \ne m$) vanish by independence ($\mathbb{E}[W_{jk} W_{jm}] = 0$):
\begin{equation}
\text{Var}(y_j) = \sum_{k=1}^{d_{\text{in}}} \mathbb{E}[W_{jk}^2] \mathbb{E}[x_k^2] = d_{\text{in}} \cdot \text{Var}(W) \cdot \text{Var}(x).
\end{equation}
\end{proof}

\subsection{Canonical Initialization Schemes}
\begin{enumerate}
    \item \textbf{LeCun Initialization (1998):} For linear and tanh activations near zero, setting $\text{Var}(W) = \frac{1}{d_{\text{in}}}$ enforces $\text{Var}(y) = \text{Var}(x)$.
    \item \textbf{Glorot / Xavier Initialization (2010):} Compromising between forward signal preservation ($\text{Var}(W) = \frac{1}{d_{\text{in}}}$) and backward gradient variance preservation ($\text{Var}(W) = \frac{1}{d_{\text{out}}}$):
    \begin{equation}
    \text{Var}(W) = \frac{2}{d_{\text{in}} + d_{\text{out}}}.
    \end{equation}
    For a continuous uniform distribution $\mathcal{U}(-r, +r)$, whose variance is $\frac{r^2}{3}$, the bounding interval is $r = \sqrt{\frac{6}{d_{\text{in}} + d_{\text{out}}}}$.
    \item \textbf{He / Kaiming Initialization (2015):} Rectified linear units (ReLU) zero out half the activations on average ($\mathbb{E}[\max(0, z)^2] = \frac{1}{2} \text{Var}(z)$). To compensate, the variance is doubled:
    \begin{equation}
    \text{Var}(W) = \frac{2}{d_{\text{in}}}, \quad r = \sqrt{\frac{6}{d_{\text{in}}}}.
    \end{equation}
\end{enumerate}

\section{Complexity Analysis}
\label{sec:complexity}

\subsection{Time Complexity}
\begin{itemize}
    \item Matrix-Matrix Product $\mathbf{X} \mathbf{W}^\top$: Computes $B \cdot d_{\text{out}}$ inner products of dimension $d_{\text{in}}$. Total multiplications: $B \cdot d_{\text{in}} \cdot d_{\text{out}}$; additions: $B \cdot (d_{\text{in}} - 1) \cdot d_{\text{out}}$.
    \item Bias Addition: $B \cdot d_{\text{out}}$ additions.
    \item Total time complexity is strictly:
    \begin{equation}
    T(B, d_{\text{in}}, d_{\text{out}}) = \mathcal{O}(B \cdot d_{\text{in}} \cdot d_{\text{out}}).
    \end{equation}
\end{itemize}

\subsection{Space Complexity}
\begin{itemize}
    \item The output matrix requires allocating $\mathbf{Y} \in \mathbb{R}^{B \times d_{\text{out}}}$.
    \item Working auxiliary memory is $\mathcal{O}(1)$.
    \begin{equation}
    S(B, d_{\text{out}}) = \mathcal{O}(B \cdot d_{\text{out}}).
    \end{equation}
\end{itemize}

\section{Worked Example}
\label{sec:worked_example}

Consider batch size $B=2$, fan-in $d_{\text{in}} = 3$, fan-out $d_{\text{out}} = 2$.
\begin{equation}
\mathbf{X} = \begin{bmatrix} 1.0 & 2.0 & -1.0 \\ 0.0 & -0.5 & 1.5 \end{bmatrix}, \quad
\mathbf{W} = \begin{bmatrix} 0.5 & -1.0 & 2.0 \\ -2.0 & 0.0 & 1.0 \end{bmatrix}, \quad
\mathbf{b} = \begin{bmatrix} 0.25 \\ -0.50 \end{bmatrix}.
\end{equation}

\paragraph{Computation for Sample 1 ($i=1$, $\mathbf{x}_1 = [1.0, 2.0, -1.0]$):}
\begin{align}
Y_{1, 1} &= (1.0)(0.5) + (2.0)(-1.0) + (-1.0)(2.0) + 0.25 = 0.5 - 2.0 - 2.0 + 0.25 = -3.25, \\
Y_{1, 2} &= (1.0)(-2.0) + (2.0)(0.0) + (-1.0)(1.0) - 0.50 = -2.0 + 0.0 - 1.0 - 0.50 = -3.50.
\end{align}

\paragraph{Computation for Sample 2 ($i=2$, $\mathbf{x}_2 = [0.0, -0.5, 1.5]$):}
\begin{align}
Y_{2, 1} &= (0.0)(0.5) + (-0.5)(-1.0) + (1.5)(2.0) + 0.25 = 0.0 + 0.5 + 3.0 + 0.25 = 3.75, \\
Y_{2, 2} &= (0.0)(-2.0) + (-0.5)(0.0) + (1.5)(1.0) - 0.50 = 0.0 + 0.0 + 1.5 - 0.50 = 1.00.
\end{align}

\begin{equation}
\mathbf{Y} = \begin{bmatrix} -3.25 & -3.50 \\ 3.75 & 1.00 \end{bmatrix}.
\end{equation}

\paragraph{Initialization Metrics (mode = \texttt{glorot\_uniform}):}
\begin{equation}
\sigma^2 = \frac{2}{3 + 2} = \frac{2}{5} = 0.40, \quad r = \sqrt{3 \times 0.40} = \sqrt{1.20} \approx 1.095445.
\end{equation}

\section{Numerical Considerations}
\label{sec:numerical}

\begin{enumerate}
    \item \textbf{Accumulation Precision:} For large fan-in dimensions ($d_{\text{in}} > 10^4$), floating-point summation order can cause drift. Double-precision scalar accumulators guarantee error within $10^{-15}$.
    \item \textbf{Cache Locality in MatMul:} While the reference CPU implementation evaluates sequentially, batched implementations transpose $\mathbf{W}$ to align with contiguous memory vectors.
\end{enumerate}

\begin{thebibliography}{9}

\bibitem{glorot2010understanding}
X.~Glorot and Y.~Bengio.
\newblock Understanding the difficulty of training deep feedforward neural networks.
\newblock In {\em Proceedings of the Thirteenth International Conference on Artificial Intelligence and Statistics}, pages 249--256, 2010.

\bibitem{he2015delving}
K.~He, X.~Zhang, S.~Ren, and J.~Sun.
\newblock Delving deep into rectifiers: Surpassing human-level performance on ImageNet classification.
\newblock In {\em Proceedings of the IEEE International Conference on Computer Vision (ICCV)}, pages 1026--1034, 2015.
\newblock DOI: \href{https://doi.org/10.1109/ICCV.2015.123}{10.1109/ICCV.2015.123}.

\bibitem{lecun1998efficient}
Y.~LeCun, L.~Bottou, G.~B.~Orr, and K.-R.~M{\"u}ller.
\newblock Efficient backprop.
\newblock In {\em Neural Networks: Tricks of the Trade}, pages 9--50. Springer, 1998.
\newblock DOI: \href{https://doi.org/10.1007/3-540-49430-8_2}{10.1007/3-540-49430-8\_2}.

\end{thebibliography}

\end{document}
